\section{Related Work}
\label{sec:related}

Related work is organized around the two alignment stages.  The first question
is where task authority comes from and how it reaches an online \actionblock.
The second is how far the resulting evidence follows that block toward physical
execution.  Section~\ref{sec:background} gives the capability-level comparison;
the discussion below explains the corresponding mechanisms and claim
boundaries.

\subsection{VLA Attacks, Safety, and Candidate Verification}

SABER uses bounded instruction edits to stress-test \vla policies
~\cite{saber2026}.  FreezeVLA studies action-freezing attacks against \vla
models, while BadRobot and RoboPAIR study jailbreaks of embodied agents and
LLM-controlled robots~\cite{freezevla2025,badrobot2025,robopair2025}.  A recent
survey organizes this broader space by multimodal attack surface, physical
effect, real-time constraint, and long-horizon propagation~\cite{vlasafety2026}.
These works establish that corrupt context can alter embodied behavior; RQ1
measures how often one such attack reaches a paired physical-risk endpoint in the
studied action-only path.

LIBERO-Safety, SafeVLA-Bench, and ForesightSafety-VLA separate task completion
from safety outcomes~\cite{liberosafety2026,safevlabench2026,
foresightsafety2026}.  SafeVLA integrates constraints during policy training,
and SAFE learns rollout-failure detection~\cite{safevla2025,safe2025}.  These
approaches modify or train the policy, whereas \system leaves the checkpoint and
policy interface unchanged and mediates its output at runtime.

SEAL samples action sequences, predicts outcomes, and selects against a
reasoning \vla's textual plan~\cite{sealvla2026}.  CoVer learns a contrastive
instruction--observation--action verifier and scales across rephrased
instructions and action candidates~\cite{covervla2026}.  Their authority comes
from a policy plan or learned candidate score.  \system instead assesses one
source block against independently anchored task state and carries that same
proposal into dispatch and effect evidence.  The approaches are complementary.
Candidate quality and execution integrity are different claims.

\subsection{Trusted Inputs and Digital-Agent Execution}

An agent-security SoK maps the broader attack-and-defense landscape
~\cite{agenticsok2026}.  StruQ separates trusted and untrusted query fields, and
SecAlign uses preference optimization for prompt-injection defense
~\cite{struq2025,secalign2025}.  Studies of deployed plugins and tool-using
agents show attacks through conversation history, tool metadata, source
ordering, and cross-tool flow~\cite{webplugins2026,toolhijacker2026,
obliinjection2026,lesdissonances2026}.  This line of work motivates separating
trusted authority from mutable model context.

CaMeL derives protected control and data flow from a trusted query, while ACE
derives trusted abstract plans and enforces data and capability barriers
~\cite{camel2026,ace2026}.  IsolateGPT, SAGA, and automated permission systems
add application isolation, agent identity, and data-access control
~\cite{isolategpt2025,saga2026,agentpermissions2026}.  AttriGuard causally
attributes tool calls, and MATE learns a policy-conditioned mobile-agent
trajectory auditor~\cite{attriguard2026,mate2026}.  Their protected objects are
digital control flows, calls, identities, permissions, or mobile-agent
trajectories.  \system applies the trusted-authority principle to a high-rate
continuous actuator transaction and follows the online source block beyond the
digital decision point.

\subsection{CPS Integrity, Robot Attestation, and Shielding}

DIAT addresses malicious modification of generated or processed data through
control-flow-backed attestation~\cite{diat2019}; CFA+ detects or prevents
control-flow hijacking~\cite{cfaplus2024}; and ARTO protects control and data
flow in real-time CPS~\cite{arto2026}.  SCAPHY correlates industrial-control
program behavior with physical behavior~\cite{scaphy2023}.  TAT addresses
robot-motion manipulation using motion-event and joint evidence
~\cite{tat2026}.  These systems begin from a protected program, execution, or
reference motion.  The protected object in \system is a newly generated
continuous block whose task-relative assessment must be connected to later
execution.

VLMPC predicts action-conditioned frames for candidate action sequences, and
SafetyChance calibrates predicted safety for image-controlled autonomy
~\cite{vlmpc2024,safetychance2024}.  Measurement-robust control barrier
functions and continuous-space shields enforce defined safety sets
~\cite{measurementrobustcbf2021,realizableshields2025}.  These methods motivate
state-conditioned qualification.  In \system, that qualification is one part of
a larger evidence lifecycle: it names both the fixed source action and the
controller configuration used to realize it, and remains limited to the
registered joint boundary.

Classic protection design requires complete mediation and fail-safe defaults
~\cite{saltzer1975}, while execution-monitoring theory characterizes properties
enforceable over event prefixes~\cite{schneider2000}.  \system applies these
principles to proposal, dispatch, receipt, effect, and phase events.  Lean checks
selected abstract transaction relations; sensing, checker coverage, and dynamics
remain explicit TCB or empirical assumptions.

Taken together, prior systems protect upstream authority, candidate selection,
downstream execution integrity, or local physical constraints.  The composition
developed here is a single lifecycle that binds independent task authority to
the exact generated block, carries the same proposal through ordered execution
evidence, and records its state-conditioned realization.
